% ------------------------------------------------------------------
%  Appendice B — Schede pratiche
%  Ricerca di riferimento: ricerca 01 §17.2; 04 §3–8; e soprattutto i
%    capitoli da cui ogni scheda è ripresa (9, 13, 14, 20, 21, 22, 23, 24)
%  Scheda nell'indice ragionato: ricerca/06_indice-ragionato.md, §6.3
%
%  STATO: prima stesura (07/10/2026), circa 3.550 parole; 19 pagine
%  (scheda B.1: 1 pagina; B.2: 2 + modulo; B.3: 1 + modulo; B.4: 3 + modulo;
%  B.5: 3; B.6: 2; B.7: 2). Testo delle schede in corpo minore (\small).
%  Chiude il registro Y-04: calendario di ripasso (cap. 9, rimando nel
%  riquadro «In pratica»), «Il metodo in una pagina» (cap. 14),
%  il semestre in 14 settimane (cap. 20, tab:semestre), la tabella
%  «Una regola, sei traduzioni» e le cinque schede per area (cap. 21,
%  tab:discipline). Aggiunge, come da indice §6.3: preparazione di uno
%  scritto e di un orale (cap. 20, con capp. 8, 13, 14, 15), protocollo
%  per chi ha la dislessia (cap. 22), lezione e prove per chi insegna
%  (capp. 23–24).
%
%  REGOLA: le schede non contengono dati, citazioni o promesse che non
%  siano già nei capitoli; riprendono, condensandoli, testi già
%  verificati. Se un capitolo cambia, va cambiata anche la sua scheda:
%  in particolare B.6 dopo la revisione esterna del cap. 22, B.5 se
%  cambia il cap. 21, B.2 se cambia la tabella 9.1.
%  Unici elementi nuovi, dichiaratamente indicativi e non presentati
%  come risultati: la riga d'esempio del calendario (giorno dopo, una
%  settimana, due, tre: coerente con la tabella 9.1 per un esame a fine
%  semestre) e la frequenza delle tre scatole (B.2).
%  Voce: «tu» lo studente; nella scheda B.7 «tu» il docente, senza
%  «noi docenti» (decisione 13 dell'indice, 07/10/2026).
%  Impaginazione: ogni scheda comincia su una pagina nuova, per poterla
%  fotocopiare da sola; i moduli da compilare sono tabelle a griglia
%  con righe vuote, su una pagina propria.
% ------------------------------------------------------------------
\chapter{Schede pratiche}
\label{app:schede}

%  Comandi usati solo in questa appendice
%  Il testo delle schede è in corpo minore, come i riquadri dei capitoli.
\newcommand{\scheda}[2]{\clearpage\normalsize\section{#1}\label{#2}\small}
\newcommand{\dacapitolo}[1]{%
  {\centering\footnotesize\itshape #1\par}\vspace{\baselineskip}}
\newcommand{\casella}{\tikz[baseline=-0.05em]\draw[line width=0.4pt]
  (0,0) rectangle (0.62em,0.62em);}
\newcommand{\vuota}{\rule{0pt}{1.45\baselineskip}}
\newcommand{\modulo}[1]{%
  \clearpage\normalsize
  {\centering{\color{rubrica}\scshape\spaziato[120]{\MakeLowercase{#1}}}\par}%
  \vspace{0.5\baselineskip}\small}
\newenvironment{controlli}
  {\begin{itemize}[label=\casella, leftmargin=1.6em, itemsep=0pt, topsep=0.2\baselineskip]}
  {\end{itemize}}
%  Le schede finiscono dove finiscono: niente spazi stirati a fondo pagina
\raggedbottom

\iniziale{Q}{ueste pagine} sono fatte per essere usate, non lette.
Raccolgono, in forma di schede, le procedure pratiche del libro: il
metodo, il calendario delle riprese, il piano del semestre, la
preparazione degli esami, le traduzioni del metodo nelle diverse
discipline, un percorso per chi ha la dislessia e, per chi insegna, la
lezione e le prove intermedie. Ogni scheda comincia su una pagina
nuova, perché possa essere fotocopiata da sola, e indica il capitolo da
cui viene.

Le schede non sostituiscono i capitoli. Dicono \emph{che cosa} fare, non
\emph{perché}: le ragioni, le prove e i loro limiti sono nei capitoli a
cui rimandano, ed è bene averli letti almeno una volta. Una procedura
di cui non si conosce il perché si abbandona alla prima difficoltà; una
di cui si conosce il perché si adatta alla propria situazione senza
perderne il cuore. Per lo stesso motivo nessuna scheda contiene dati o
promesse che non siano già nel testo: dove il capitolo dice che una
pratica è provata solo in parte, o soltanto proposta, lo dice anche la
scheda.

Un consiglio d'uso. Non serve fotocopiarle tutte. Scegline una o due --
per cominciare, il metodo in una pagina e il calendario di ripasso -- e
tienile dove studi: sulla scrivania, nel quaderno, come prima pagina
dell'agenda. I moduli da compilare (il calendario delle riprese, il
piano del semestre, il quaderno degli errori) si ricopiano ogni volta
che servono; se ti sembrano piccoli, stampali dal PDF gratuito del
libro ingrandendoli a tutta pagina. La licenza del libro permette di
copiarli e di distribuirli senza scopo di lucro, citandone la fonte:
chi insegna può darli ai propri studenti.

%% =================================================================
\scheda{Il metodo in una pagina}{sch:metodo}
\dacapitolo{Dal capitolo~\ref{cap:metodo}. Le ragioni: capitoli~\ref{cap:recupero}--\ref{cap:insegnare}.}

{\small
\begin{enumerate}[leftmargin=1.4em, itemsep=0.25\baselineskip]
  \item \textbf{Prima} \emph{(cinque minuti)}. Guarda l'argomento
        nuovo -- titoli, obiettivi, domande finali -- e prova a
        rispondere per iscritto a due o tre domande, anche se non sai
        la risposta. Poi rispondi a memoria a una domanda
        sull'argomento precedente.
  \item \textbf{Durante}. Capisci, non trascrivere. Annota le idee
        principali, gli esempi e le domande che ti vengono. Alla fine
        di ogni paragrafo chiediti: che cosa ho letto? perché è vero?
        come si collega a ciò che so? Negli esempi svolti, spiegati
        ogni passaggio.
  \item \textbf{Subito dopo} \emph{(venti minuti, lo stesso giorno)}.
        A libro chiuso, scrivi tutto ciò che ricordi; poi controlla
        sugli appunti e correggi con un altro colore. Prepara da cinque
        a dieci domande sui punti essenziali.
  \item \textbf{Il giorno dopo} \emph{(un quarto d'ora per lezione)}.
        Rispondi alle domande senza guardare. Ristudia subito solo ciò
        che hai sbagliato.
  \item \textbf{Ogni settimana} \emph{(da mezz'ora a un'ora per
        insegnamento, sempre lo stesso giorno)}. Rispondi alle domande
        di \emph{tutte} le settimane, mescolate, fino al criterio;
        esercizi misti; una spiegazione ad alta voce, senza appunti.
  \item \textbf{Prima dell'esame} \emph{(le ultime due settimane)}.
        Niente materiale nuovo. Simulazioni della prova, nelle sue
        condizioni. Ripresa di tutto il programma, mescolato. La notte
        prima si dorme.
  \item \textbf{Sempre}. Telefono lontano, pause vere, sonno
        sufficiente. E una risposta scritta alla domanda \enquote{a
        che cosa mi serve questo, e per chi?}.
\end{enumerate}}

\vspace{0.5\baselineskip}
{\centering\small\itshape Rileggere ti dà la sensazione di sapere;\\
rispondere ti dice se sai.\par}

%% =================================================================
\scheda{Il calendario di ripasso}{sch:calendario}
\dacapitolo{Dal capitolo~\ref{cap:tempo}; per il ripasso a successive
riprese, anche il capitolo~\ref{cap:metodo}.}

\subsection*{Quanto aspettare tra una ripresa e l'altra}

\begin{center}\small
  \begin{tabular}{@{}>{\raggedright\arraybackslash}p{0.6\textwidth}>{\raggedright\arraybackslash}p{0.34\textwidth}@{}}
    \toprule
    \textit{Devi ricordare per\dots} & \textit{Intervallo tra le riprese} \\
    \midrule
    una settimana (una prova in itinere) & 1--2 giorni \\
    un mese & circa una settimana \\
    tre mesi (un esame a fine semestre) & 2--3 settimane \\
    un anno (un concorso, un esame di Stato) & 2--4 settimane \\
    per sempre & ogni pochi mesi \\
    \bottomrule
  \end{tabular}
\end{center}
{\footnotesize Valori orientativi (tabella~\ref{tab:intervalli}). Nel
dubbio, meglio un intervallo un po' più lungo che uno più breve.\par}

\subsection*{Le regole della ripresa}

\begin{enumerate}[leftmargin=1.4em]
  \item \emph{In ogni sessione}: domande a libro chiuso, fino al
        criterio. Se sbagli, controlli la risposta e rimetti la domanda
        in fondo al mazzo, finché ogni domanda non ha avuto almeno una
        risposta giusta.
  \item \emph{Da una sessione all'altra}: ciò che sai bene torna più di
        rado; ciò che sbagli torna presto. Una domanda esce dal mazzo
        solo quando hai risposto bene in più sessioni, in giorni
        diversi.
  \item \emph{Sul calendario}: ogni ripresa è un appuntamento scritto,
        con un giorno e un'ora, non un proposito.
  \item \emph{Senza scoraggiarti}: la prima sessione è lunga, le
        successive sono brevi. Un calendario ragionevole e costante vale
        più di uno perfetto e abbandonato.
\end{enumerate}

\subsection*{Tre situazioni tipiche}

{\small\emph{Una prova in itinere tra una settimana.} Primo giorno: studia per
capire, fai il primo recupero a libro chiuso, prepara le domande.
Secondo giorno: domande fino al criterio, ristudia gli errori. Quarto o
quinto giorno: di nuovo. Sesto giorno: una simulazione della prova.
Settimo: una breve ripresa, e a letto presto.

\emph{Un esame a fine semestre.} Per ogni insegnamento, una sessione
settimanale di ripresa, sempre lo stesso giorno, con le domande della
settimana e di quelle precedenti. Ogni argomento torna così cinque o
sei volte prima della sessione (scheda~\ref{sch:semestre}).

\emph{Un concorso o un esame di Stato tra un anno.} Dividi il programma
in blocchi e studiali uno dopo l'altro; ma dedica ogni settimana una
parte fissa del tempo, per esempio un quarto, a riprendere i blocchi
già studiati, a distanza di due, tre o quattro settimane.\par}

\subsection*{Il mazzo e le scatole}

{\small Se usi schede di carta, dividile in scatole: chi risponde bene
sposta la scheda nella scatola successiva, che si riprende più di rado;
chi sbaglia la riporta nella prima. Tre scatole bastano: per esempio, la
prima a ogni sessione, la seconda ogni settimana o due, la terza una
volta al mese. Un'applicazione di flashcard fa lo stesso calcolo da
sola, purché tu risponda prima di girare la scheda e valuti
onestamente la risposta.\par}

%  ---------------------------------------------------------- modulo
\modulo{Calendario delle riprese}

{\small\noindent Insegnamento \hrulefill\par
\vspace{0.4\baselineskip}
\noindent Esame il \hrulefill\quad Intervallo scelto \hrulefill\par}
\vspace{0.6\baselineskip}

{\footnotesize\noindent In ogni casella scrivi la data prevista per la
ripresa; quando l'hai fatta, segnala con una crocetta. La prima riga è
un esempio per un esame a fine semestre: il giorno dopo, poi dopo una,
due e tre settimane.\par}
\vspace{0.6\baselineskip}

\begingroup
\setlength{\tabcolsep}{3pt}\setlength{\arrayrulewidth}{0.3pt}
\begin{center}\footnotesize
  \begin{tabular}{|>{\raggedright\arraybackslash}p{0.315\textwidth}|*{5}{>{\centering\arraybackslash}p{0.1\textwidth}|}}
    \hline
    \rule{0pt}{1.1\baselineskip}\textit{Argomento} & \textit{Studio} & \textit{1ª} & \textit{2ª} & \textit{3ª} & \textit{4ª} \\
    \hline
    \vuota\textit{Media e mediana} & \textit{6/10} & \textit{7/10} & \textit{14/10} & \textit{28/10} & \textit{18/11} \\
    \hline
    \vuota & & & & & \\ \hline
    \vuota & & & & & \\ \hline
    \vuota & & & & & \\ \hline
    \vuota & & & & & \\ \hline
    \vuota & & & & & \\ \hline
    \vuota & & & & & \\ \hline
    \vuota & & & & & \\ \hline
    \vuota & & & & & \\ \hline
    \vuota & & & & & \\ \hline
    \vuota & & & & & \\ \hline
    \vuota & & & & & \\ \hline
    \vuota & & & & & \\ \hline
    \vuota & & & & & \\ \hline
  \end{tabular}
\end{center}
\endgroup

%% =================================================================
\scheda{Il semestre in quattordici settimane}{sch:semestre}
\dacapitolo{Dal capitolo~\ref{cap:semestre} (tabella~\ref{tab:semestre});
il lavoro di ogni settimana è nella scheda~\ref{sch:metodo}. Le
settimane sono indicative.}

\begingroup
\setlength{\tabcolsep}{4pt}
\begin{center}\footnotesize
  \begin{tabular}{@{}>{\raggedright\arraybackslash}p{0.17\textwidth}>{\raggedright\arraybackslash}p{0.68\textwidth}c@{}}
    \toprule
    \textit{Quando} & \textit{Che cosa fare} & \\
    \midrule
    Settimana 1 & Il piano del semestre (modulo alla pagina seguente);
      un quaderno, o un mazzo, di domande per insegnamento. & \casella \\
    \addlinespace
    Settimane 2--4 & Il metodo, ogni giorno e ogni settimana. Alla fine
      della quarta, una prova a libro chiuso su tutto ciò che è stato
      spiegato. & \casella \\
    \addlinespace
    Settimane 5--6 & Le domande che sai bene tornano più di rado, quelle
      sbagliate subito. Esercizi misti dove l'esame è di problemi; una
      sessione a coppie alla settimana dove l'esame è orale. & \casella \\
    \addlinespace
    Settimane 7--8 & Prove in itinere, dove ci sono: una simulazione
      pochi giorni prima; dopo la prova, l'analisi degli errori. Le
      domande della prova entrano nel mazzo. & \casella \\
    \addlinespace
    Settimane 9--12 & Nelle sessioni settimanali, una parte fissa per le
      domande delle prime settimane. Alla decima, una mappa di ogni
      programma fatta a memoria e corretta sul libro: ciò che manca è la
      lista delle lacune. & \casella \\
    \addlinespace
    Settimane 13--14 & Prima lettura del manuale completata. Calendario
      della sessione: quali giorni a quale esame, con le simulazioni già
      fissate. & \casella \\
    \addlinespace
    Sessione & Niente materiale nuovo. Riprese mescolate di tutto il
      programma, simulazioni degli scritti e degli orali, sonno. Gli
      esami del semestre nel primo appello utile. & \casella \\
    \bottomrule
  \end{tabular}
\end{center}
\endgroup



%  ---------------------------------------------------------- modulo
\modulo{Il piano del semestre}

{\small\noindent Semestre \hrulefill\quad compilato il \hrulefill\par}
\vspace{0.6\baselineskip}

{\footnotesize\noindent Da compilare nella prima settimana di lezione:
richiede un'ora. Se hai più di tre insegnamenti, usa una seconda
copia.\par}
\vspace{0.6\baselineskip}

\begingroup
\setlength{\tabcolsep}{3pt}\setlength{\arrayrulewidth}{0.3pt}
\begin{center}\footnotesize
  \begin{tabular}{|>{\raggedright\arraybackslash}p{0.27\textwidth}|*{3}{>{\raggedright\arraybackslash}p{0.2\textwidth}|}}
    \hline
    \rule{0pt}{1.55\baselineskip}Insegnamento & & & \\ \hline
    \rule{0pt}{1.6\baselineskip}Crediti & & & \\ \hline
    \rule{0pt}{1.55\baselineskip}Esame: scritto, orale o entrambi & & & \\ \hline
    \rule{0pt}{1.55\baselineskip}Prove in itinere (date) & & & \\ \hline
    \rule{0pt}{1.55\baselineskip}Appello scelto & & & \\ \hline
    \rule{0pt}{1.55\baselineskip}Giorno e ora della sessione settimanale & & & \\ \hline
    \rule{0pt}{1.55\baselineskip}Materiale per allenarsi (testi d'esame, esercizi) & & & \\ \hline
  \end{tabular}
\end{center}
\endgroup

\vspace{0.4\baselineskip}
{\small\noindent Quali esami di questo semestre voglio chiudere nel
primo appello, e che cosa devo fare ogni settimana perché accada?\par
\vspace{0.9\baselineskip}\noindent\hrulefill\par
\vspace{0.6\baselineskip}
\noindent In quali settimane si accumulano le scadenze?\par
\vspace{0.9\baselineskip}\noindent\hrulefill\par}

%% =================================================================
\scheda{Prepararsi a uno scritto e a un orale}{sch:esami}
\dacapitolo{Dal capitolo~\ref{cap:semestre}; le ragioni nei
capitoli~\ref{cap:recupero}, \ref{cap:alternanza}, \ref{cap:insegnare}
e~\ref{cap:corpo}.}

\subsection*{Dopo ogni lezione, in venti minuti}

\begin{enumerate}[leftmargin=1.4em]
  \item \emph{A libro chiuso} (dieci minuti): scrivi tutto ciò che
        ricordi della lezione.
  \item \emph{Confronta} (cinque minuti): completa con un altro colore
        ciò che mancava; sono le tue prime lacune.
  \item \emph{Domande} (cinque minuti): da cinque a dieci, nella colonna
        di sinistra degli appunti, di almeno due tipi diversi.
\end{enumerate}

\begin{center}\small
  \begin{tabular}{@{}>{\raggedright\arraybackslash}p{0.27\textwidth}>{\raggedright\arraybackslash}p{0.65\textwidth}@{}}
    \toprule
    \textit{Tipo di domanda} & \textit{Un esempio di Statistica} \\
    \midrule
    di memoria & Che cos'è la mediana? \\
    di spiegazione & Perché la mediana non risente dei valori estremi? \\
    di applicazione & Per descrivere i redditi di un Paese, media o
      mediana? Perché? \\
    di collegamento & Con una lunga coda verso destra, la media è
      maggiore o minore della mediana? \\
    \bottomrule
  \end{tabular}
\end{center}

\subsection*{Per uno scritto}

\begin{controlli}
  \item Testi d'esame passati, o esercizi dello stesso tipo, usati
        \emph{alla fine} dello studio, non all'inizio.
  \item Svolti nelle condizioni dell'esame: tempo contato, niente
        appunti, niente telefono, niente soluzioni a portata di mano.
  \item Problemi mescolati: prima di risolverne uno, scrivi di che tipo
        è.
  \item Correzione severa, e ogni errore nel quaderno degli errori, con
        la sua specie.
\end{controlli}

\subsection*{Il quaderno degli errori}

\begin{center}\small
  \begin{tabular}{@{}>{\raggedright\arraybackslash}p{0.22\textwidth}>{\raggedright\arraybackslash}p{0.33\textwidth}>{\raggedright\arraybackslash}p{0.35\textwidth}@{}}
    \toprule
    \textit{Specie} & \textit{Che cosa è successo} & \textit{Rimedio} \\
    \midrule
    Conoscenza (C) & non sapevi la definizione, la formula, il teorema &
      ristudiarlo e metterlo nel mazzo delle riprese \\
    \addlinespace
    Riconoscimento (R) & sapevi, ma non hai capito quale strumento usare &
      esercizi misti, con il tipo di problema scritto prima \\
    \addlinespace
    Esecuzione (E) & un segno, un passaggio saltato, il tempo finito &
      simulazioni a tempo e una revisione finale sistematica \\
    \bottomrule
  \end{tabular}
\end{center}

\subsection*{Per un orale}

\begin{controlli}
  \item Per ogni argomento: esponilo ad alta voce, senza guardare --
        una definizione, un esempio, un collegamento. Dove ti inceppi,
        lì c'è il lavoro.
  \item Registrati e riascoltati: dove sei vago, dove ripeti le stesse
        formule, dove perdi il filo?
  \item Fatti interrogare da un compagno, con domande scelte da lui, che
        saltino da una parte all'altra del programma e chiedano
        collegamenti.
  \item Le \enquote{domande che il professore fa sempre}: per
        interrogarti, non per imparare a memoria le risposte.
\end{controlli}

\begin{riquadro}{Una sessione a coppie in trenta minuti}
Prima, ciascuno studia lo stesso argomento e prepara cinque domande,
di memoria e di applicazione. \emph{Dieci minuti}: A spiega a B, a
libro chiuso, la prima metà; B interrompe con \enquote{perché?} e
\enquote{fammi un esempio}, e alla fine controlla sul testo che cosa
manca. \emph{Dieci minuti}: si invertono i ruoli sulla seconda metà.
\emph{Dieci minuti}: ciascuno pone all'altro le sue cinque domande; le
risposte si controllano insieme. Alla fine, ciascuno annota le lacune
da riprendere (capitolo~\ref{cap:insegnare}).
\end{riquadro}

\subsection*{Gli ultimi giorni, e il giorno dell'esame}

\begin{controlli}
  \item Niente materiale nuovo; riprese mescolate di tutto il programma,
        con più spazio per ciò che sbagli.
  \item La notte prima si dorme. La mattina, se vuoi, un breve recupero
        sulle domande principali; non rileggere tutto il manuale.
  \item Se all'orale arriva il vuoto: chiedi di riformulare la domanda,
        o comincia da ciò che ricordi con sicurezza -- un esempio, una
        definizione, un collegamento. Spesso il resto torna mentre
        parli.
\end{controlli}

\begin{riquadro}{Se sei già in ritardo}
\begin{enumerate}[leftmargin=1.4em]
  \item \emph{Fai subito il punto}: per ogni argomento del programma,
        scrivi a memoria ciò che ne sai. In un'ora saprai che cosa è a
        posto, che cosa è debole, che cosa manca.
  \item \emph{Distribuisci anche il poco tempo che hai}: ogni giorno un
        po' di tutto, non un argomento al giorno fino alla vigilia.
  \item \emph{Studia interrogandoti}: dopo la prima lettura, niente
        riletture, solo domande.
  \item \emph{Non rinunciare al sonno}.
  \item \emph{Valuta con lucidità} l'appello successivo, se il tempo
        davvero non basta (capitolo~\ref{cap:metodo}).
\end{enumerate}
\end{riquadro}

%  ---------------------------------------------------------- modulo
\modulo{Quaderno degli errori}

{\small\noindent Insegnamento \hrulefill\par}
\vspace{0.6\baselineskip}

{\footnotesize\noindent Specie: C = conoscenza, R = riconoscimento,
E = esecuzione. Nella colonna \enquote{Ripreso} segna la data in cui hai
rifatto un esercizio dello stesso tipo, senza errori.\par}
\vspace{0.6\baselineskip}

\begingroup
\setlength{\tabcolsep}{3pt}\setlength{\arrayrulewidth}{0.3pt}
\begin{center}\footnotesize
  \begin{tabular}{|>{\centering\arraybackslash}p{0.1\textwidth}|>{\raggedright\arraybackslash}p{0.25\textwidth}|>{\centering\arraybackslash}p{0.09\textwidth}|>{\raggedright\arraybackslash}p{0.27\textwidth}|>{\centering\arraybackslash}p{0.11\textwidth}|}
    \hline
    \rule{0pt}{1.1\baselineskip}\textit{Data} & \textit{Esercizio} & \textit{Specie} & \textit{Rimedio} & \textit{Ripreso} \\
    \hline
    \vuota & & & & \\ \hline
    \vuota & & & & \\ \hline
    \vuota & & & & \\ \hline
    \vuota & & & & \\ \hline
    \vuota & & & & \\ \hline
    \vuota & & & & \\ \hline
    \vuota & & & & \\ \hline
    \vuota & & & & \\ \hline
    \vuota & & & & \\ \hline
    \vuota & & & & \\ \hline
    \vuota & & & & \\ \hline
    \vuota & & & & \\ \hline
    \vuota & & & & \\ \hline
    \vuota & & & & \\ \hline
  \end{tabular}
\end{center}
\endgroup

%% =================================================================
\scheda{Ogni disciplina ha il suo studio}{sch:discipline}
\dacapitolo{Dal capitolo~\ref{cap:discipline} (tabella~\ref{tab:discipline}),
dove sono descritti gli studi su cui si regge la colonna \enquote{quanto
è provato}.}

La regola è la stessa per tutte le discipline: richiamare, distanziare,
alternare, spiegare. Cambia la forma del recupero, e cambia quanto ne
sappiamo.

\begin{center}\small
  \begin{tabular}{@{}>{\raggedright\arraybackslash}p{0.22\textwidth}>{\raggedright\arraybackslash}p{0.42\textwidth}>{\raggedright\arraybackslash}p{0.26\textwidth}@{}}
    \toprule
    \textit{Area} & \textit{Recuperare vuol dire\dots} & \textit{Quanto è provato} \\
    \midrule
    Matematica, fisica, ingegneria & risolvere dal foglio bianco;
      riconoscere il tipo di problema in una serie mista & bene, anche
      in un corso universitario \\
    \addlinespace
    Medicina e scienze della vita & domande con correzione, schede
      distanziate, casi mescolati & bene, ma su campioni piccoli; le
      app solo osservate \\
    \addlinespace
    Giurisprudenza ed economia & richiamare la regola e applicarla a un
      caso nuovo; spiegare ad alta voce & poco: correlazioni, studi
      americani \\
    \addlinespace
    Discipline umanistiche & rispondere per iscritto a «perché?»;
      interrogare la fonte & per la storia, a scuola; per il resto,
      principi generali \\
    \addlinespace
    Lingue moderne & produrre la parola e la frase & bene per il
      lessico \\
    \addlinespace
    Lingue classiche & produrre le forme; tradurre e ritradurre &
      nessuna prova specifica \\
    \bottomrule
  \end{tabular}
\end{center}

\subsection*{Matematica, fisica, ingegneria}

\begin{enumerate}[leftmargin=1.4em, itemsep=0pt, topsep=0.2\baselineskip]
  \item \emph{All'inizio di un argomento}: due o tre esempi svolti,
        spiegandoti ogni passaggio; poi copri la soluzione e rifalli.
  \item \emph{Poi}: problemi dal foglio bianco. Guarda la soluzione solo
        dopo un tentativo vero, e annota nel quaderno degli errori di
        che specie era l'errore (scheda~\ref{sch:esami}).
  \item \emph{Ogni settimana}: una serie mista, senza titoli, con
        problemi del capitolo nuovo e di quelli precedenti. Prima di
        risolvere ciascun problema, scrivi di che tipo è.
  \item \emph{Teoremi e formule}: enunciali a memoria e, quando puoi,
        ricavali; una formula che sai ricavare non la dimentichi.
\end{enumerate}

\subsection*{Medicina e scienze della vita}

\begin{enumerate}[leftmargin=1.4em, itemsep=0pt, topsep=0.2\baselineskip]
  \item \emph{Ogni giorno}: poche schede nuove, scritte da te a partire
        dalla lezione, e prima di quelle le riprese in scadenza. Mai
        schede nuove se le riprese sono indietro.
  \item \emph{Domande di più tipi}: nomi e definizioni, ma anche
        \enquote{perché} e \enquote{che cosa succede se}; risposte
        brevi scritte o dette, non solo scelta multipla.
  \item \emph{Casi e immagini mescolati}: tracciati, preparati, casi
        clinici di apparati e diagnosi diverse nella stessa sessione;
        prova la diagnosi prima di leggere la soluzione.
  \item \emph{Collega}: per ogni struttura, malattia o farmaco, una
        domanda che tenga insieme anatomia, fisiologia e clinica.
\end{enumerate}

\subsection*{Giurisprudenza ed economia}

\begin{enumerate}[leftmargin=1.4em, itemsep=0pt, topsep=0.2\baselineskip]
  \item \emph{Per ogni istituto o concetto}: definizione a memoria,
        norma o modello di riferimento, un caso in cui si applica e uno
        in cui non si applica.
  \item \emph{Casi mescolati}: brevi fatti, presi dal manuale, dalle
        sentenze o inventati da un compagno, da qualificare: quale
        istituto, quale regola, perché.
  \item \emph{Il codice come prova}: leggi, chiudi, riformula, confronta.
  \item \emph{Modelli economici}: grafici e relazioni ricavati dal
        foglio bianco, con le domande \enquote{che cosa succede se}.
  \item \emph{Parla}: ogni settimana, spiega ad alta voce un argomento a
        qualcuno che ti faccia domande.
\end{enumerate}

\subsection*{Discipline umanistiche}

\begin{enumerate}[leftmargin=1.4em, itemsep=0pt, topsep=0.2\baselineskip]
  \item \emph{Dopo ogni capitolo o testo}: due risposte scritte a libro
        chiuso, a un \enquote{perché?} e a un \enquote{che cosa vuol
        dire?}.
  \item \emph{Davanti a una fonte}: chi scrive, per chi, quando, perché;
        con quale altra fonte la confronto.
  \item \emph{Cronologie, schemi di un'opera, argomenti}: ricostruiti a
        memoria e corretti sul libro.
  \item \emph{Confronti}: due autori, correnti o epoche nella stessa
        sessione, con la domanda \enquote{in che cosa differiscono?};
        e, quando puoi, discutine ad alta voce con qualcuno.
\end{enumerate}

\subsection*{Lingue straniere e classiche}

\begin{enumerate}[leftmargin=1.4em, itemsep=0pt, topsep=0.2\baselineskip]
  \item \emph{Lessico}: schede in entrambi i versi, soprattutto dal
        significato alla parola; riprese a intervalli crescenti, per
        mesi, non solo fino all'esame.
  \item \emph{Grammatica}: forme che si confondono, mescolate nella
        stessa sessione; per ogni frase, quale forma e perché.
  \item \emph{Lingue moderne}: ogni settimana parla, ascolta e scrivi;
        le schede non bastano.
  \item \emph{Lingue classiche}: paradigmi prodotti a memoria; testi
        ripresi più volte; tradurre, e dopo qualche giorno ritradurre
        senza guardare l'originale.
\end{enumerate}

%% =================================================================
\scheda{Studiare con la dislessia}{sch:dislessia}
\dacapitolo{Dal capitolo~\ref{cap:dsa}. Non è una scheda clinica né
legale: per le misure che ti spettano valgono la legge, le Linee guida
e il servizio del tuo ateneo.}

\subsection*{Un capitolo di manuale, in sette passi}

\begin{enumerate}[leftmargin=1.4em, itemsep=0pt, topsep=0.2\baselineskip]
  \item \emph{Accesso}: ascolta il capitolo con la sintesi vocale,
        seguendo con gli occhi titoli, figure e parole in grassetto.
  \item \emph{Primo recupero}: subito dopo, di' a voce -- o registra --
        tutto ciò che ricordi. Riascolta solo le parti mancanti.
  \item \emph{Mappa}: costruisci una mappa essenziale a partire da ciò
        che hai detto, non dal libro.
  \item \emph{Il giorno dopo}: copri la mappa e ricostruiscila a voce;
        poi controlla.
  \item \emph{Nei giorni seguenti}: brevi sessioni di domande orali,
        da solo, con un compagno o con un assistente che ti interroghi.
  \item \emph{Prima dell'esame}: simulazioni orali, e scritti di prova
        con il tempo che ti spetta.
  \item \emph{Sempre}: sessioni brevi e distanziate, pause vere,
        sonno.
\end{enumerate}

\subsection*{Che cosa, secondo le prove, non aiuta}

Gli allenamenti computerizzati della memoria di lavoro (migliorano i
compiti allenati, non la lettura né lo studio); i caratteri tipografici
\enquote{per dislessici}; le lenti e i filtri colorati; studiare secondo
un presunto \enquote{stile di apprendimento} (capitoli~\ref{cap:dsa}
e~\ref{cap:miti}).

\subsection*{Compensare l'accesso, non il pensiero}

Uno strumento è prezioso quando toglie un ostacolo strumentale, come la
lettura lenta o la grafia faticosa; diventa un danno quando toglie le
attività che producono apprendimento: richiamare, elaborare,
riprendere nel tempo. È un criterio \emph{proposto}, non un risultato
sperimentale.

\begin{center}\footnotesize
  \begin{tabular}{@{}>{\raggedright\arraybackslash}p{0.24\textwidth}>{\raggedright\arraybackslash}p{0.33\textwidth}>{\raggedright\arraybackslash}p{0.33\textwidth}@{}}
    \toprule
    \textit{Strumento} & \textit{Aiuta se\dots} & \textit{Ostacola se\dots} \\
    \midrule
    Sintesi vocale & dopo l'ascolto richiami a voce & riascolti senza
      richiamare \\
    \addlinespace
    Mappe e schemi & li costruisci tu e ti interroghi coprendoli & usi
      mappe fatte da altri \\
    \addlinespace
    Registrazioni & colmi una lacuna precisa & riascolti lezioni intere \\
    \addlinespace
    Calcolatrice, formulari & liberano attenzione per ragionare & rinunci
      a imparare ciò che potresti imparare \\
    \addlinespace
    Tempo aggiuntivo & toglie all'esame il peso della lentezza & riduce
      la pratica durante lo studio \\
    \addlinespace
    Intelligenza artificiale & legge, ti interroga, commenta ciò che hai
      prodotto & riassume e risponde al posto tuo \\
    \bottomrule
  \end{tabular}
\end{center}

{\footnotesize Che cosa ti spetta è una questione di diritto, e non è
un invito a rinunciare a una misura prevista. Come conviene studiare,
invece, vale per tutti: se all'esame potrai usare una mappa, costruiscila
a memoria e riprendila nei giorni precedenti, perché sia l'indice di ciò
che sai, non il testo da leggere.\par}

\subsection*{Cinque passi con il servizio di ateneo}

\begin{controlli}
  \item \emph{All'inizio dell'anno}, non alla vigilia del primo esame:
        consegna la certificazione al servizio e chiedi un colloquio.
  \item \emph{Per ogni insegnamento}, all'inizio del corso o almeno un
        mese prima dell'esame, concorda con il docente le misure che ti
        servono, spiegando perché.
  \item \emph{Per i test d'ingresso}, presenta domanda e certificazione
        entro i termini del bando.
  \item \emph{Usa il tutorato}, quando c'è: è un servizio previsto, non
        un favore.
  \item \emph{Se una misura ti viene negata}, o se i docenti si
        comportano in modo molto diverso, torna al servizio e al
        delegato del rettore.
\end{controlli}

%% =================================================================
\scheda{Per chi insegna: la lezione e le prove}{sch:docenti}
\dacapitolo{Dai capitoli~\ref{cap:docenti} e~\ref{cap:valutare}; le
ragioni anche nei capitoli~\ref{cap:recupero} e~\ref{cap:tempo}.}

\subsection*{Prima che il corso cominci}

\begin{controlli}
  \item Le date delle prove in itinere, fissate e comunicate alla prima
        lezione.
  \item I testi degli esami passati e qualche domanda d'esempio per
        l'orale, a disposizione degli studenti.
  \item Una descrizione chiara di come si svolge l'esame.
  \item Un archivio di domande da costruire settimana dopo settimana, da
        cui ripescare anche quelle delle lezioni precedenti.
  \item Alla prima lezione, qualche minuto per spiegare il perché di
        quiz e prove frequenti, con i dati del
        capitolo~\ref{cap:recupero}.
\end{controlli}

\subsection*{La lezione in cinque momenti (in due ore, due o tre blocchi)}

\begin{enumerate}[leftmargin=1.4em, itemsep=0pt, topsep=0.2\baselineskip]
  \item \emph{Ripasso iniziale} (5 minuti): tre domande, a cui tutti
        rispondono per iscritto a libro chiuso, sulla lezione
        precedente e su una di qualche settimana prima; poi le
        risposte corrette, subito.
  \item \emph{Spiegazione a blocchi brevi} (15--20 minuti ciascuno):
        un'idea per blocco, con un esempio svolto; niente slide
        affollate, niente testo da leggere mentre parli.
  \item \emph{Una domanda a tutta l'aula} dopo ogni blocco: concettuale,
        non \enquote{ci sono domande?}; ciascuno risponde da solo, poi
        si confronta con il vicino, poi la discussione.
  \item \emph{Un problema a coppie} (10 minuti) su ciò che si è appena
        spiegato; tu giri tra i banchi e ascolti gli errori.
  \item \emph{Domande di uscita} (3 minuti): \enquote{Scrivi ciò che
        ricordi della lezione di oggi} e \enquote{Che cosa non ti è
        chiaro?}. Le risposte ti dicono da dove ripartire.
\end{enumerate}

\subsection*{Da un esame unico a quattro momenti di recupero (un corso di dodici settimane)}

\begin{enumerate}[leftmargin=1.4em, itemsep=0pt, topsep=0.2\baselineskip]
  \item \emph{Ogni settimana}: un quiz breve (cinque domande, dieci
        minuti, in aula o in rete), \emph{dopo} la spiegazione, sugli
        argomenti della settimana e su uno delle settimane precedenti;
        correzione subito; peso piccolo sul voto, o un bonus.
  \item \emph{Sesta settimana}: una prima prova in itinere sulle prime
        cinque settimane, con un riscontro scritto entro dieci giorni.
  \item \emph{Undicesima settimana}: una seconda prova, cumulativa: due
        terzi delle domande sugli argomenti nuovi, un terzo su quelli
        della prima parte.
  \item \emph{Esame finale}: su tutto il programma, con un peso ridotto
        per chi ha svolto le prove intermedie, ma senza parti
        \enquote{già fatte}.
\end{enumerate}

\subsection*{Un riscontro che aiuta}

\begin{controlli}
  \item Risponde a tre domande: \emph{dove sto andando?}, \emph{come sto
        andando?}, \emph{che cosa faccio adesso?}
  \item Parla del compito e del processo (\enquote{hai applicato la
        regola senza controllare le condizioni}), non della persona
        (\enquote{sei bravo}).
  \item Arriva tra una prova e l'altra, non dopo l'esame finale.
  \item Quando gli studenti sono troppi, è collettivo: dieci minuti di
        lezione sugli errori più frequenti della prova in itinere.
  \item La visione del compito diventa una spiegazione degli errori,
        non una contrattazione sul voto.
\end{controlli}

\subsection*{Tre cose da evitare}

Le prove in itinere che \enquote{liberano} una parte del programma:
dicono agli studenti che quella parte si può dimenticare; meglio prove
che alleggeriscono l'esame senza toglierne nulla. Le prove con i
problemi sempre in ordine di capitolo: all'esame vero bisogna
riconoscerli. E un nuovo modo di insegnare con gli esami di sempre: gli
studenti continueranno a studiare per gli esami di sempre.

\normalsize
\finecapitolo
\flushbottom
